% GENERATED FILE. Edit canonical lessons, then run npm run generate:books.
% canonical-source-hash: fnv1a64:9f221e38fe8dea73
% canonical-lessons: PA-C307-majjh, PA-C307-katta, PA-C307-chona, PA-C307-than, PA-C307-charna

\chapter{The Buffalo}
\label{ch:pa-a2-the-buffalo}
\hlchaptermodality{307}

\begin{chapteropening}
\textbf{By the end of this chapter:} \emph{I can talk about the buffalo and its milk.}

Complete the last lesson of chapter 307: I can talk about the buffalo and its milk.
\end{chapteropening}

\section[\texorpdfstring{\pa{ਮੱਝ} (majjh)}{ਮੱਝ (majjh)}]{\pa{ਮੱਝ} (majjh) — a buffalo}

\label{lesson:PA-C307-majjh}



\begin{quote}
Before the new one: say the Punjabi for rice straw, then the Punjabi for a sack.
\end{quote}

\subsection*{You'll want to know: \pa{ਮੱਝ}}
\textbf{\pa{ਮੱਝ}} — \emph{majjh} — \textquotedblleft{}a buffalo\textquotedblright{}.

\begin{quote}
\pa{ਸਾਡੀ} \pa{ਮੱਝ} \pa{ਬਹੁਤ} \pa{ਦੁੱਧ} \pa{ਦਿੰਦੀ} \pa{ਹੈ।}
\end{quote}

\begin{quote}
\emph{sāḍī majjh bahut dudh dindī hai.}
\end{quote}

\begin{quote}
Our buffalo gives a lot of milk.
\end{quote}

A feminine noun.

\begin{cousinweb}
\textbf{\pa{ਮੱਝ}} (\emph{majjh}) is inherited from a reconstructed \emph{*mahyā}, \textquotedblleft{}buffalo\textquotedblright{}; CDIAL lists Punjabi \emph{majjh} under it (CDIAL 9980).
\end{cousinweb}

\begin{grammarlens}[title={What you've built}]
One more word for this chapter.
\end{grammarlens}

\subsection*{Your turn}
\begin{itemize}
\raggedright
  \item \emph{Say it:} \emph{majjh}
  \item \emph{Say it:} \emph{majjh}, once more
  \item \emph{Say it:} \emph{sāḍī majjh bahut dudh dindī hai.} — Our buffalo gives a lot of milk.
  \item \emph{Recall:} say the Punjabi for rice straw, then the Punjabi for a sack, then say \emph{majjh} again
\end{itemize}

\begin{tcolorbox}[breakable,colback=teal!4,colframe=teal!35!black,title={Before you move on}]
What does \pa{ਮੱਝ} mean? (\textquotedblleft{}A buffalo\textquotedblright{}.) Say it once more.
\end{tcolorbox}
\section[\texorpdfstring{\pa{ਕੱਟਾ} (kaṭṭā)}{ਕੱਟਾ (kaṭṭā)}]{\pa{ਕੱਟਾ} (kaṭṭā) — a buffalo calf}

\label{lesson:PA-C307-katta}



\begin{quote}
Before the new one: say the Punjabi for a sack, then the Punjabi for a buffalo.
\end{quote}

\subsection*{You'll want to know: \pa{ਕੱਟਾ}}
\textbf{\pa{ਕੱਟਾ}} — \emph{kaṭṭā} — \textquotedblleft{}a buffalo calf\textquotedblright{}.

\begin{quote}
\pa{ਮੱਝ} \pa{ਦਾ} \pa{ਕੱਟਾ} \pa{ਛੋਟਾ} \pa{ਹੈ।}
\end{quote}

\begin{quote}
\emph{majjh dā kaṭṭā chhoṭā hai.}
\end{quote}

\begin{quote}
The buffalo's calf is small.
\end{quote}

A masculine noun.

\begin{cousinweb}
\textbf{\pa{ਕੱਟਾ}} (\emph{kaṭṭā}) is inherited from a reconstructed \emph{*kaṭṭa}, \textquotedblleft{}young male animal\textquotedblright{}; CDIAL lists Punjabi \emph{kaṭṭā}, \textquotedblleft{}yearling buffalo\textquotedblright{}, under it (CDIAL 2645).
\end{cousinweb}

\begin{grammarlens}[title={What you've built}]
One more word for this chapter.
\end{grammarlens}

\subsection*{Your turn}
\begin{itemize}
\raggedright
  \item \emph{Say it:} \emph{kaṭṭā}
  \item \emph{Say it:} \emph{kaṭṭā}, once more
  \item \emph{Say it:} \emph{majjh dā kaṭṭā chhoṭā hai.} — The buffalo's calf is small.
  \item \emph{Recall:} say the Punjabi for a sack, then the Punjabi for a buffalo, then say \emph{kaṭṭā} again
  \item \emph{Write it:} copy \textbf{\pa{ਕੱਟਾ}} once, checking the whole word
\end{itemize}

\begin{tcolorbox}[breakable,colback=teal!4,colframe=teal!35!black,title={Before you move on}]
What does \pa{ਕੱਟਾ} mean? (\textquotedblleft{}A buffalo calf\textquotedblright{}.) Say it once more.
\end{tcolorbox}
\section[\texorpdfstring{\pa{ਚੋਣਾ} (choṇā)}{ਚੋਣਾ (choṇā)}]{\pa{ਚੋਣਾ} (choṇā) — to milk an animal}

\label{lesson:PA-C307-chona}



\begin{quote}
Before the new one: say the Punjabi for a buffalo, then the Punjabi for a buffalo calf.
\end{quote}

\subsection*{You'll want to know: \pa{ਚੋਣਾ}}
\textbf{\pa{ਚੋਣਾ}} — \emph{choṇā} — \textquotedblleft{}to milk an animal\textquotedblright{}.

\begin{quote}
\pa{ਮੈਂ} \pa{ਸਵੇਰੇ} \pa{ਮੱਝ} \pa{ਚੋਵਾਂਗਾ।}
\end{quote}

\begin{quote}
\emph{maiṁ savere majjh chovāṅgā.}
\end{quote}

\begin{quote}
I will milk the buffalo in the morning.
\end{quote}

Past (I …): \pa{ਮੈਂ} \pa{ਮੱਝ} \pa{ਚੋਈ}, \emph{maiṁ majjh choī}.

\begin{cousinweb}
\textbf{\pa{ਚੋਣਾ}} (\emph{choṇā}) is inherited from Sanskrit \emph{cyōtati}, \textquotedblleft{}flows, trickles\textquotedblright{}; CDIAL lists Punjabi \emph{coṇā}, \textquotedblleft{}to leak, milk\textquotedblright{}, under it (CDIAL 4948).
\end{cousinweb}

\begin{grammarlens}[title={What you've built}]
One more word for this chapter.
\end{grammarlens}

\subsection*{Your turn}
\begin{itemize}
\raggedright
  \item \emph{Say it:} \emph{choṇā}
  \item \emph{Say it:} \emph{choṇā}, once more
  \item \emph{Say it:} \emph{maiṁ savere majjh chovāṅgā.} — I will milk the buffalo in the morning.
  \item \emph{Recall:} say the Punjabi for a buffalo, then the Punjabi for a buffalo calf, then say \emph{choṇā} again
  \item \emph{Write it:} copy \textbf{\pa{ਚੋਣਾ}} once, checking the whole word
\end{itemize}

\begin{tcolorbox}[breakable,colback=teal!4,colframe=teal!35!black,title={Before you move on}]
What does \pa{ਚੋਣਾ} mean? (\textquotedblleft{}To milk an animal\textquotedblright{}.) Say it once more.
\end{tcolorbox}
\section[\texorpdfstring{\pa{ਥਣ} (thaṇ)}{ਥਣ (thaṇ)}]{\pa{ਥਣ} (thaṇ) — a teat; an udder}

\label{lesson:PA-C307-than}



\begin{quote}
Before the new one: say the Punjabi for a buffalo calf, then the Punjabi for to milk an animal.
\end{quote}

\subsection*{You'll want to know: \pa{ਥਣ}}
\textbf{\pa{ਥਣ}} — \emph{thaṇ} — \textquotedblleft{}a teat; an udder\textquotedblright{}.

\begin{quote}
\pa{ਮੱਝ} \pa{ਦੇ} \pa{ਥਣਾਂ} \pa{ਵਿੱਚ} \pa{ਬਹੁਤ} \pa{ਦੁੱਧ} \pa{ਹੈ।}
\end{quote}

\begin{quote}
\emph{majjh de thaṇāṁ vich bahut dudh hai.}
\end{quote}

\begin{quote}
There is a lot of milk in the buffalo's teats.
\end{quote}

A masculine noun.

\begin{cousinweb}
\textbf{\pa{ਥਣ}} (\emph{thaṇ}) is inherited from Sanskrit \emph{stana}, \textquotedblleft{}female breast, udder\textquotedblright{}; CDIAL lists Punjabi \emph{thaṇ}, \textquotedblleft{}teat\textquotedblright{}, under it (CDIAL 13666).
\end{cousinweb}

\begin{grammarlens}[title={What you've built}]
One more word for this chapter.
\end{grammarlens}

\subsection*{Your turn}
\begin{itemize}
\raggedright
  \item \emph{Say it:} \emph{thaṇ}
  \item \emph{Say it:} \emph{thaṇ}, once more
  \item \emph{Say it:} \emph{majjh de thaṇāṁ vich bahut dudh hai.} — There is a lot of milk in the buffalo's teats.
  \item \emph{Recall:} say the Punjabi for a buffalo calf, then the Punjabi for to milk an animal, then say \emph{thaṇ} again
  \item \emph{Write it:} copy \textbf{\pa{ਥਣ}} once, checking the whole word
\end{itemize}

\begin{tcolorbox}[breakable,colback=teal!4,colframe=teal!35!black,title={Before you move on}]
What does \pa{ਥਣ} mean? (\textquotedblleft{}A teat; an udder\textquotedblright{}.) Say it once more.
\end{tcolorbox}
\section[\texorpdfstring{\pa{ਚਰਨਾ} (charnā)}{ਚਰਨਾ (charnā)}]{\pa{ਚਰਨਾ} (charnā) — to graze}

\label{lesson:PA-C307-charna}



\begin{quote}
Before the new one: say the Punjabi for to milk an animal, then the Punjabi for a teat.
\end{quote}

\subsection*{You'll want to know: \pa{ਚਰਨਾ}}
\textbf{\pa{ਚਰਨਾ}} — \emph{charnā} — \textquotedblleft{}to graze\textquotedblright{}.

\begin{quote}
\pa{ਗਾਂ} \pa{ਦਰਿਆ} \pa{ਦੇ} \pa{ਕਿਨਾਰੇ} \pa{ਚਰਦੀ} \pa{ਹੈ।}
\end{quote}

\begin{quote}
\emph{gāṁ dariā de kināre chardī hai.}
\end{quote}

\begin{quote}
The cow grazes on the bank of the river.
\end{quote}

Past (it …): \pa{ਉਹ} \pa{ਬੱਕਰੀ} \pa{ਚਰੀ}, \emph{oh bakkrī charī}.

\begin{cousinweb}
\textbf{\pa{ਚਰਨਾ}} (\emph{charnā}) is inherited from Sanskrit \emph{carati}, \textquotedblleft{}moves, grazes\textquotedblright{}; CDIAL lists Punjabi \emph{carnā}, \textquotedblleft{}to graze\textquotedblright{}, under it (CDIAL 4686).
\end{cousinweb}

\begin{grammarlens}[title={What you've built}]
One more word for this chapter.
\end{grammarlens}

\subsection*{Your turn}
\begin{itemize}
\raggedright
  \item \emph{Say it:} \emph{charnā}
  \item \emph{Say it:} \emph{charnā}, once more
  \item \emph{Say it:} \emph{gāṁ dariā de kināre chardī hai.} — The cow grazes on the bank of the river.
  \item \emph{Recall:} say the Punjabi for to milk an animal, then the Punjabi for a teat, then say \emph{charnā} again
  \item \emph{Write it:} copy \textbf{\pa{ਚਰਨਾ}} once, checking the whole word
\end{itemize}

\begin{tcolorbox}[breakable,colback=teal!4,colframe=teal!35!black,title={Before you move on}]
What does \pa{ਚਰਨਾ} mean? (\textquotedblleft{}To graze\textquotedblright{}.) Say it once more.
\end{tcolorbox}
